\documentclass[../main.tex]{subfiles}

\begin{document}

\ifSubfilesClassLoaded{

    \tableofcontents
    
}{}

\chapter{ Smooth Maps }

\section{Smooth Functions and Smooth Maps}

\subsection{Smooth Maps Between Manifolds}

\begin{definition}{Smooth Maps Between Manifolds}{}
    Let $M, N$ be smooth manifolds, and let $F: M \to N$ be any map. We say that $F$ is a \dfntxt{smooth map} if for every $p \in M$, there exist smooth charts $(U, \varphi)$ containing $p$ and $(V, \psi)$ containing $F(p)$ such that $F(U) \subseteq V$ and the composite map $\psi \circ F \circ \varphi^{-1}$ is smooth from $\varphi(U)$ to $\psi(V)$ .
\end{definition}

\begin{remark}
    \begin{itemize}
        \item By taking \(  N= V= \mathbb{R} ^{k}  \) and \(  \psi =\mathrm{Id}:\mathbb{R} ^{k}\to \mathbb{R} ^{k}  \), we get the definition of smoothness of real-valued or vector-valed functions.  
    \end{itemize}
    
\end{remark}

\end{document}